\documentclass[12pt]{article}

% Packages for math, formatting, and layout
\usepackage{amsmath}
\usepackage{amssymb}
\usepackage{geometry}
\usepackage{hyperref}
\usepackage{enumitem}

% Set page margins
\geometry{a4paper, margin=1in}

\title{\textbf{Dynamic Programming Algorithms}}
\author{}
\date{}

\begin{document}

\maketitle


\subsection*{1) Policy Evaluation}
Policy evaluation is the process of calculating the state-value function ($V^\pi(s)$). It tells us the goodness or badness of a certain policy ($\pi$). For this, we use the Bellman Expectation Equation for the state-value function:
$$V_{k+1}(s) = \sum_a \pi(a|s) \sum_{s', r} p(s', r | s, a) \left[ r + \gamma V_k(s') \right]$$

\noindent\textbf{Pseudocode:}
\begin{itemize}[label=$\bullet$]
    \item \textbf{Input:}
    \begin{itemize}[label=$\circ$]
        \item The policy $\pi$ which we want to evaluate.
        \item An estimation threshold $\theta > 0$ to determine the accuracy of our estimation.
        \item Transition probabilities / Transition Matrix.
        \item Discount factor $\gamma \in [0,1]$.
    \end{itemize}
    \item \textbf{Process:}
    \begin{itemize}[label=$\circ$]
        \item Set the value of every state arbitrarily.
        \item Set the value of the terminal state to 0.
        \item Set $\Delta$ to 0.
        \item \textbf{Loop} until $\Delta < \theta$:
        \begin{itemize}[label=--]
            \item For each state in the state space:
            \item Store the current value of the state in a variable.
            \item Update the value of the state using the Bellman Expectation Equation.
            \item Check how much the state's value changed.
            \item Set $\Delta \leftarrow \max(\text{current } \Delta, \text{absolute difference between the old and new state value})$.
        \end{itemize}
    \end{itemize}
    \item \textbf{Output:} The final value of each state, which is now a close approximation of the value function $V^\pi(s)$ for the given policy.
\end{itemize}

\subsection*{2) Policy Iteration}
Policy Iteration is an algorithm used to discover the optimal policy.

\newpage
\noindent\textbf{Pseudocode:}
\begin{itemize}[label=$\bullet$]
    \item \textbf{Input:}
    \begin{itemize}[label=$\circ$]
        \item An estimation threshold $\theta > 0$ to determine the accuracy of our estimation.
        \item Transition probabilities / Transition Matrix.
        \item Discount factor $\gamma \in [0,1]$.
    \end{itemize}
    \item \textbf{Process:}
    \begin{itemize}[label=$\circ$]
        \item Set the value of every state arbitrarily.
        \item Set $\Delta$ to 0.
        \item Start with an arbitrary policy and perform \textbf{Policy Evaluation}.
        \item Then perform \textbf{Policy Improvement}, for each state:
        \begin{itemize}[label=--]
            \item Store the current action and calculate the expected return for all possible actions.
            \item Select the action with the highest expected return.
            \item Update the policy to choose this action.
            \item If the selected action differs from the previous action, mark the policy as unstable.
        \end{itemize}
        \item Check whether the policy changed: If no action changed in any state, stop; otherwise, repeat Policy Evaluation with the improved policy.
    \end{itemize}
    \item \textbf{Output:} Optimal policy $\pi^*$ and the optimal value function $V^*(s)$ corresponding to that policy.
\end{itemize}

\subsection*{3) Value Iteration}
Value Iteration is an algorithm which helps in finding the optimal strategy faster. For this, we use the Bellman Optimality Equation of the state-value function:
$$V_{k+1}(s) = \max_a \sum_{s', r} p(s', r \mid s, a) \left[ r + \gamma V_k(s') \right]$$

\noindent\textbf{Pseudocode:}
\begin{itemize}[label=$\bullet$]
    \item \textbf{Input:}
    \begin{itemize}[label=$\circ$]
        \item An estimation threshold $\theta > 0$ to determine the accuracy of our estimation.
        \item Transition probabilities / Transition Matrix.
        \item Discount factor $\gamma \in [0,1]$.
    \end{itemize}
    \item \textbf{Process:}
    \begin{itemize}[label=$\circ$]
        \item Set the value of every state arbitrarily.
        \item Set the value of the terminal state to 0.
        \item \textbf{Loop} until the maximum change in state values ($\Delta$) becomes smaller than $\theta$:
        \begin{itemize}[label=--]
            \item Set $\Delta$ to 0.
            \item Repeatedly update the value of each state using the Bellman Optimality Equation.
        \end{itemize}
        \item Extract the optimal policy by choosing, in each state, the action that yields the highest expected return according to the converged value function.
    \end{itemize}
    \item \textbf{Output:} Optimal policy $\pi^*$ and the optimal value function $V^*(s)$ corresponding to that policy.
\end{itemize}

\vspace{0.5cm}

\section*{Comparison between Monte Carlo (MC) and Temporal Difference (TD)}
We will take an example of a small toy MDP for this comparison.

\subsection*{Toy MDP Setup}
Imagine a $2 \times 2$ grid:
\begin{itemize}[label=$\bullet$]
    \item Top Left (TL) $\rightarrow$ Top Right (TR)
    \item Bottom Left (BL) $\rightarrow$ Bottom Right (BR)
\end{itemize}
\textbf{BR is our Goal.}

\begin{itemize}[label=$\bullet$]
    \item \textbf{Actions:} Up, Down, Left, Right.
    \item \textbf{Rewards:} Every step taken costs $-1$ \& stepping into goal (BR) gives $+10$.
    \item \textbf{Discount factor ($\gamma$):} $1$ (i.e., no discounting).
    \item \textbf{Learning Rate ($\alpha$):} $0.5$.
\end{itemize}

\noindent Initially, $V(TL)=0, V(TR)=0, V(BL)=0, V(BR)=0$.

\vspace{0.3cm}
\noindent Let's simulate one specific episode starting at BL:
$$ \text{BL} \xrightarrow{\text{UP}} \text{TL (} R=-1 \text{)} \xrightarrow{\text{RIGHT}} \text{TR (} R=-1 \text{)} \xrightarrow{\text{DOWN}} \text{BR (} R=+10 \text{)} $$

\noindent Now compare between Monte Carlo and Temporal Difference:

\subsection*{1) Monte Carlo (MC) Execution}
MC must wait until it hits a terminal state as it calculates cumulative return ($G_t$) from each state.

\vspace{0.3cm}
\noindent\textbf{Work Backwards:}
\begin{itemize}[label=$\bullet$]
    \item From TR $\rightarrow$ only one step and total reward is 10. $\therefore G_{TR} = 10$.
    \item From TL $\rightarrow G_{TL} = (-1) + 10 = 9$.
    \item From BL $\rightarrow G_{BL} = (-1) + (-1) + 10 = 8$.
\end{itemize}

\noindent\textbf{Now update the values:}
\begin{itemize}[label=$\bullet$]
    \item \textbf{Update TR} $\rightarrow V(TR) \leftarrow V(TR) + \alpha [G_{TR} - V(TR)]$ \\
    $V(TR) \leftarrow 0 + 0.5[10 - 0] = 5$
    \item \textbf{Update TL} $\rightarrow V(TL) \leftarrow V(TL) + \alpha [G_{TL} - V(TL)]$ \\
    $V(TL) \leftarrow 0 + 0.5[9 - 0] = 4.5$
    \item \textbf{Update BL} $\rightarrow V(BL) \leftarrow V(BL) + \alpha [G_{BL} - V(BL)]$ \\
    $V(BL) \leftarrow 0 + 0.5[8 - 0] = 4$
\end{itemize}

\noindent \textbf{We got:} $V(TR)=5, \ V(TL)=4.5, \ V(BL)=4$.

\subsection*{2) Temporal Difference (TD) Execution}
TD updates the value immediately after every move.

\begin{itemize}[label=$\bullet$]
    \item \textbf{First step (BL $\rightarrow$ TL)} $\rightarrow R + \gamma V(TL) = -1 + 1(0) = -1$ \\
    $V(BL) \leftarrow V(BL) + 0.5[-1 - 0] = -0.5$
    \item \textbf{Second step (TL $\rightarrow$ TR)} $\rightarrow R + \gamma V(TR) = -1 + 1(0) = -1$ \\
    $V(TL) \leftarrow V(TL) + 0.5[-1 - 0] = -0.5$
    \item \textbf{Third step (TR $\rightarrow$ BR)} $\rightarrow R + \gamma V(BR) = 10 + 1(0) = 10$ \\
    $V(TR) \leftarrow V(TR) + 0.5[10 - 0] = 5$
\end{itemize}

\noindent \textbf{We got:} $V(TR)=5, \ V(TL)=-0.5, \ V(BL)=-0.5$.

\subsection*{Value Comparison}
\begin{table}[h!]
\centering
\renewcommand{\arraystretch}{1.5}
\begin{tabular}{|c|c|c|}
\hline
& \textbf{MC} & \textbf{TD} \\ \hline
\textbf{$V(TR)$} & $5$ & $5$ \\ \hline
\textbf{$V(TL)$} & $4.5$ & $-0.5$ \\ \hline
\textbf{$V(BL)$} & $4$ & $-0.5$ \\ \hline
\end{tabular}
\end{table}

\end{document}